\section{Webová prezentace fotbalového klubu s databází zápasů a správou
obsahu ve
WordPressu}\label{webovuxe1-prezentace-fotbalovuxe9ho-klubu-s-databuxe1zuxed-zuxe1pasux16f-a-spruxe1vou-obsahu-ve-wordpressu}

\begin{center}\rule{0.5\linewidth}{0.5pt}\end{center}

\textbf{Maturitní práce}

\textbf{Autor:} Lukáš Bejček \textbf{Obor:} Informační technologie
\textbf{Školní rok:} 2025/2026

\begin{center}\rule{0.5\linewidth}{0.5pt}\end{center}

\newpage

\subsection{Prohlášení}\label{prohluxe1ux161enuxed}

Prohlašuji, že jsem maturitní práci na téma \emph{Webová prezentace
fotbalového klubu s databází zápasů a správou obsahu ve WordPressu}
vypracoval samostatně pod vedením vedoucího práce a s použitím odborné
literatury a dalších informačních zdrojů, které jsou citovány v práci a
uvedeny v seznamu použité literatury.

V Mýtě dne \ldots\ldots\ldots\ldots\ldots\ldots.. Podpis:
\ldots\ldots\ldots\ldots\ldots\ldots\ldots.

\newpage

\subsection{Poděkování}\label{podux11bkovuxe1nuxed}

Děkuji vedoucímu práce za odborné vedení, cenné rady a připomínky při
zpracování maturitní práce.

\newpage

\subsection{Abstrakt}\label{abstrakt}

Tato maturitní práce se zabývá návrhem a implementací webové prezentace
fotbalového klubu TJ Slavoj Mýto. Cílem bylo vytvořit moderní,
responzivní web, který umožňuje správci klubu spravovat zápasy, týmy,
hráče, galerie, sponzory a kontakty bez znalosti programování. Web je
postaven na redakčním systému WordPress s vlastním tématem založeným na
frameworku Bootstrap 5 a vlastním pluginem pro rozšíření administrace.
Místo externích placených pluginů byla zvolena vlastní implementace
klíčových funkcí -- registrace typů obsahu, vlastních polí, filtrů a
uživatelských rolí -- přímo v kódu tématu a pluginu. Práce popisuje
analýzu původního řešení, návrh datového modelu, implementaci šablon a
administračních nástrojů, testování, příručku administrátora a postup
nasazení na hosting.

\textbf{Klíčová slova:} WordPress, fotbalový klub, webová prezentace,
custom post types, Bootstrap 5, PHP, responzivní design

\newpage

\subsection{Obsah}\label{obsah}

\emph{(Obsah je generován automaticky při exportu do PDF -- viz
\texttt{docs/generate-pdf.sh})}

\newpage

\subsection{1 Úvod}\label{uxfavod}

Fotbalový klub TJ Slavoj Mýto je amatérský sportovní klub působící v
obci Mýto u Rokycan. Klub má několik týmů od mužů po mládežnické
kategorie a potřebuje moderní webovou prezentaci, na které budou
přehledně zobrazeny zápasy, týmy, hráčské soupisky, fotogalerie,
partneři klubu a kontakty na vedení.

Původní web klubu byl charakteristický nepřehledným uspořádáním obsahu a
graficky zastaralým vzhledem. Cílem této práce je vytvořit graficky
přijatelný a přehledný web s plně funkčním a spravovatelným WordPressem,
kde veškerý obsah lze spravovat přes standardní administrační rozhraní.

Práce je členěna do následujících kapitol. Kapitola 2 analyzuje původní
řešení a porovnává možné přístupy k implementaci. Kapitola 3 popisuje
implementaci -- datový model, šablony, administrační nástroje a
bezpečnost. Kapitola 4 dokumentuje testování a ověření funkčnosti.
Kapitola 5 slouží jako příručka administrátora pro správu obsahu webu.
Kapitola 6 popisuje nasazení na produkční hosting. Kapitola 7 obsahuje
závěr s návrhy na rozšíření.

Celý projekt je dostupný jako open-source repozitář na GitHubu a je
připraven k lokálnímu nasazení pomocí XAMPP.

\begin{center}\rule{0.5\linewidth}{0.5pt}\end{center}

\subsection{2 Analýza problému}\label{analuxfdza-probluxe9mu}

\subsubsection{2.1 Původní stav webu}\label{pux16fvodnuxed-stav-webu}

Původní web TJ Slavoj Mýto (http://slavojmyto.cz/) byl charakteristický
nepřehledným uspořádáním obsahu a graficky zastaralým vzhledem. Stránky
obsahují sekce: úvodní stránku, historii klubu, výbor (kontakty), týmy
se soupiskami, kalendář zápasů po sezónách, fotogalerie a sponzory.

Ověřením živého webu byly zjištěny tyto nedostatky:

\begin{enumerate}
\def\labelenumi{\arabic{enumi}.}
\item
  \textbf{Nepřehledná navigace} -- struktura menu neodpovídá očekávání
  návštěvníka. Zápasy jsou skryté pod položkou „Kalendář'' (nikoliv
  „Zápasy''), kontaktní údaje na vedení klubu pod „Výbor'' v sekci „O
  klubu'' (nikoliv „Kontakty''). Chybí sekce aktualit/novinek.
\item
  \textbf{Nepřehledný výpis zápasů} -- zápasy jsou zobrazeny jako dlouhý
  seznam seskupený podle týmů na jedné stránce za celou sezónu. Chybí
  jakékoliv filtrování (podle týmu, stavu zápasu) a stránkování. Pro
  nalezení konkrétního zápasu je nutné procházet celou stránku.
\item
  \textbf{Graficky zastaralý vzhled} -- vizuální styl webu neodpovídá
  moderním standardům webového designu. Layout se přizpůsobuje obrazovce
  pouze omezeně a celkový grafický dojem působí zastarale.
\item
  \textbf{Chybějící správa obsahu a funkce} -- web postrádá standardní
  administrační rozhraní pro správu obsahu, vyhledávání, detaily
  jednotlivých zápasů, profily hráčů a aktuality. Galerie používá
  dvouúrovňový systém (tým → ročník → album), ale bez filtrování.
\end{enumerate}

Z původního webu bylo možné převzít vizuální identitu (logo, barvy
klubu), strukturu sekcí a dostupná data (soupisky hráčů, výsledky
zápasů, kontakty na výbor, sponzory). Technická a datová vrstva musela
být navržena a implementována od základu v redakčním systému WordPress,
který zajistí plnohodnotnou správu veškerého obsahu přes standardní
administrační rozhraní.

\subsubsection{2.2 Požadavky na nové
řešení}\label{poux17eadavky-na-novuxe9-ux159eux161enuxed}

Na základě zadání maturitní práce {[}1{]} a potřeb klubu byly stanoveny
tyto požadavky:

\begin{itemize}
\tightlist
\item
  Evidence zápasů s výsledky, střelci a filtrováním podle sezóny, týmu a
  stavu
\item
  Evidence týmů se soupiskami hráčů a realizačním týmem
\item
  Fotogalerie členěná podle týmů a sezón
\item
  Kontakty na vedení klubu s možností řazení
\item
  Přehled sponzorů s logy a odkazy
\item
  Stránka historie klubu
\item
  Správa veškerého obsahu přes administraci WordPress bez zásahu do kódu
\item
  Responzivní design (desktop i mobilní zařízení)
\item
  Minimální JavaScript -- preferovat serverové řešení v PHP
\end{itemize}

\subsubsection{2.3 Volba technologií}\label{volba-technologiuxed}

\paragraph{2.3.1 CMS -- WordPress}\label{cms-wordpress}

Pro implementaci byl zvolen redakční systém WordPress {[}2{]},
nejrozšířenější CMS na světě s podílem přes 40 \% všech webových
stránek. WordPress nabízí:

\begin{itemize}
\tightlist
\item
  Vestavěné administrační rozhraní pro správu obsahu
\item
  Systém custom post types a taxonomií pro strukturovaná data
\item
  Šablonovou hierarchii pro mapování URL na šablony
\item
  REST API pro případné budoucí rozšíření
\item
  Rozsáhlou dokumentaci a komunitu {[}3{]}
\end{itemize}

\paragraph{2.3.2 CSS framework -- Bootstrap
5}\label{css-framework-bootstrap-5}

Pro responzivní layout byl zvolen framework Bootstrap 5 {[}4{]}, který
poskytuje:

\begin{itemize}
\tightlist
\item
  Grid systém pro responzivní rozvržení stránky
\item
  Utility třídy pro rychlé stylování (spacing, barvy, typography)
\item
  Hotové komponenty (navbar, cards, badges, forms)
\item
  Mobile-first přístup k responzivnímu designu
\end{itemize}

Bootstrap je načítán z CDN {[}5{]}, čímž odpadá nutnost build nástrojů a
zároveň se využívá cache prohlížeče.

\paragraph{2.3.3 Lokální vývojové prostředí --
XAMPP}\label{lokuxe1lnuxed-vuxfdvojovuxe9-prostux159eduxed-xampp}

Pro lokální vývoj a testování byl zvolen balík XAMPP {[}6{]}, který
obsahuje Apache webserver, MySQL databázi a PHP interpret v jedné
instalaci. XAMPP je dostupný zdarma pro Windows, macOS i Linux.

\paragraph{2.3.4 Grafický návrh --
Figma}\label{grafickuxfd-nuxe1vrh-figma}

Grafický návrh webu byl vytvořen v nástroji Figma {[}7{]} ve variantách
pro desktop i mobilní zařízení.

\subsubsection{2.4 Vlastní implementace vs.~externí
pluginy}\label{vlastnuxed-implementace-vs.-externuxed-pluginy}

Zadání práce doporučuje řadu externích pluginů: Custom Post Type UI
{[}8{]} pro registraci CPT, Advanced Custom Fields {[}9{]} pro správu
vlastních polí, FacetWP pro filtrování obsahu, The Events Calendar pro
kalendář zápasů, Modula pro galerie a User Role Editor pro správu rolí.

Po analýze jsem zvolili \textbf{vlastní implementaci} klíčových funkcí
přímo v kódu tématu a vlastním pluginu. Důvody:

{\def\LTcaptype{none} % do not increment counter
\begin{longtable}[]{@{}
  >{\raggedright\arraybackslash}p{(\linewidth - 4\tabcolsep) * \real{0.3333}}
  >{\raggedright\arraybackslash}p{(\linewidth - 4\tabcolsep) * \real{0.3333}}
  >{\raggedright\arraybackslash}p{(\linewidth - 4\tabcolsep) * \real{0.3333}}@{}}
\toprule\noalign{}
\begin{minipage}[b]{\linewidth}\raggedright
Aspekt
\end{minipage} & \begin{minipage}[b]{\linewidth}\raggedright
Externí pluginy
\end{minipage} & \begin{minipage}[b]{\linewidth}\raggedright
Vlastní implementace
\end{minipage} \\
\midrule\noalign{}
\endhead
\bottomrule\noalign{}
\endlastfoot
Kontrola nad kódem & Omezená -- závislost na cizím kódu & Plná -- kód je
čitelný a obhajitelný \\
Závislosti & Více pluginů = více bodů selhání & Pouze WordPress core +
vlastní kód \\
Výkon & Každý plugin přidává zátěž & Minimální overhead \\
Údržba & Nutnost aktualizovat pluginy & Kód je pod vlastní správou \\
Licence & Některé pluginy jsou placené & Vše zdarma (GPL-2.0) \\
\end{longtable}
}

Jediným nainstalovaným pluginem je \textbf{Slavoj Custom Fields} --
vlastní plugin vytvořený pro tento projekt, který rozšiřuje
administrační prostředí.

\subsubsection{2.5 Custom post types vs.~vlastní databázové
tabulky}\label{custom-post-types-vs.-vlastnuxed-databuxe1zovuxe9-tabulky}

WordPress nabízí dva přístupy k ukládání strukturovaných dat: custom
post types (CPT) s meta poli, nebo vlastní databázové tabulky.

Pro tento projekt byly zvoleny CPT z následujících důvodů {[}3{]}:

\begin{enumerate}
\def\labelenumi{\arabic{enumi}.}
\tightlist
\item
  \textbf{Integrace s WordPress admin} -- CPT automaticky získají
  administrační rozhraní (přehled, editace, mazání, koš) bez nutnosti
  psát vlastní CRUD logiku.
\item
  \textbf{Podpora taxonomií} -- termy (sezóna, kategorie týmu) lze
  přiřazovat nativně.
\item
  \textbf{WP\_Query} -- dotazování probíhá standardním WordPress API,
  které řeší caching, stránkování i bezpečnost.
\item
  \textbf{REST API} -- parametr
  \texttt{show\_in\_rest\ =\textgreater{}\ true} zpřístupní data přes
  JSON API.
\item
  \textbf{Šablonová hierarchie} -- WordPress automaticky mapuje URL na
  šablony (\texttt{/zapasy/} → \texttt{archive-zapas.php}).
\end{enumerate}

Alternativa s vlastními tabulkami by vyžadovala ruční implementaci
administrace, validace, stránkování a URL routingu.

\subsubsection{2.6 Archivní šablony vs.~stránkové
šablony}\label{archivnuxed-ux161ablony-vs.-struxe1nkovuxe9-ux161ablony}

Pro zobrazení seznamů CPT (zápasy, týmy, galerie) byly zvoleny
\textbf{archivní šablony} (\texttt{archive-zapas.php}) místo stránkových
šablon (\texttt{page-zapasy.php}). Důvody {[}3{]}:

\begin{itemize}
\tightlist
\item
  WordPress automaticky vytvoří čistou URL \texttt{/zapasy/} bez
  nutnosti zakládat stránku v administraci
\item
  Stránkování funguje nativně
\item
  Taxonomické filtry se integrují přirozeně přes GET parametry
\end{itemize}

Stránkové šablony zůstávají pouze pro statický obsah: historie,
kontakty, sponzoři.

\begin{center}\rule{0.5\linewidth}{0.5pt}\end{center}

\subsection{3 Implementace}\label{implementace}

\subsubsection{3.1 Architektura
řešení}\label{architektura-ux159eux161enuxed}

Projekt se skládá ze dvou komponent:

\textbf{Téma \texttt{tj-slavoj-myto}} obsahuje prezentační logiku a
registraci datového modelu:

\begin{itemize}
\tightlist
\item
  registrace custom post types a taxonomií
\item
  šablony pro front-end (homepage, archivy, detaily)
\item
  CSS styly rozdělené do komponentových souborů
\item
  pomocné funkce (řazení, vazby, počítání hráčů)
\item
  globální filtry pro konzistentní chování taxonomií
\end{itemize}

\textbf{Plugin \texttt{slavoj-custom-fields}} poskytuje administrační
nástroje:

\begin{itemize}
\tightlist
\item
  vlastní sloupce v admin přehledech (datum, tým, skóre, číslo dresu)
\item
  dropdown filtry pro filtrování v administraci
\item
  řaditelné sloupce (datum zápasu, číslo dresu)
\item
  nástrojová stránka pro správu a seedování ukázkových dat
\end{itemize}

Toto rozdělení odpovídá doporučené WordPress architektuře {[}3{]}: téma
řeší „jak vypadá'', plugin řeší „jak se spravuje''.

\subsubsection{3.2 Adresářová
struktura}\label{adresuxe1ux159ovuxe1-struktura}

\begin{verbatim}
web/
└── wp-content/
    ├── themes/tj-slavoj-myto/
    │   ├── functions.php              # Registrace CPT, taxonomií, helperů
    │   ├── front-page.php             # Homepage
    │   ├── archive-zapas.php          # Seznam zápasů s filtry
    │   ├── archive-tym.php            # Seznam týmů
    │   ├── archive-galerie.php        # Přehled fotoalb
    │   ├── single-zapas.php           # Detail zápasu
    │   ├── single-tym.php             # Profil týmu se soupiskou
    │   ├── single-hrac.php            # Profil hráče
    │   ├── single-galerie.php         # Fotoalbum s lightboxem
    │   ├── page-historie.php          # Historie klubu
    │   ├── page-kontakty.php          # Kontakty výboru
    │   ├── page-sponzori.php          # Partneři klubu
    │   ├── template-parts/            # Znovupoužitelné komponenty
    │   └── assets/css/                # Stylové soubory
    └── plugins/slavoj-custom-fields/
        └── slavoj-custom-fields.php   # Administrační plugin
\end{verbatim}

\subsubsection{3.3 Datový model}\label{datovuxfd-model}

\paragraph{3.3.1 Custom post types}\label{custom-post-types}

Projekt definuje šest vlastních typů obsahu registrovaných v
\texttt{functions.php} pomocí funkce \texttt{register\_post\_type()}
{[}3{]}:

{\def\LTcaptype{none} % do not increment counter
\begin{longtable}[]{@{}
  >{\raggedright\arraybackslash}p{(\linewidth - 6\tabcolsep) * \real{0.2500}}
  >{\raggedright\arraybackslash}p{(\linewidth - 6\tabcolsep) * \real{0.2500}}
  >{\raggedright\arraybackslash}p{(\linewidth - 6\tabcolsep) * \real{0.2500}}
  >{\raggedright\arraybackslash}p{(\linewidth - 6\tabcolsep) * \real{0.2500}}@{}}
\toprule\noalign{}
\begin{minipage}[b]{\linewidth}\raggedright
CPT
\end{minipage} & \begin{minipage}[b]{\linewidth}\raggedright
Slug
\end{minipage} & \begin{minipage}[b]{\linewidth}\raggedright
Meta pole
\end{minipage} & \begin{minipage}[b]{\linewidth}\raggedright
Účel
\end{minipage} \\
\midrule\noalign{}
\endhead
\bottomrule\noalign{}
\endlastfoot
\textbf{Zápasy} & \texttt{zapas} & \texttt{datum\_zapasu},
\texttt{cas\_zapasu}, \texttt{domaci}, \texttt{hoste}, \texttt{skore},
\texttt{strelci} & Evidence zápasů \\
\textbf{Týmy} & \texttt{tym} & \texttt{hlavni\_trener},
\texttt{asistent\_trenera}, \texttt{zdravotnik}, \texttt{tym\_slug} &
Profily týmů \\
\textbf{Hráči} & \texttt{hrac} & \texttt{cislo}, \texttt{rok\_narozeni},
\texttt{tym\_slug} & Soupisky hráčů \\
\textbf{Galerie} & \texttt{galerie} & --- & Fotoalba \\
\textbf{Sponzoři} & \texttt{sponzor} & \texttt{web\_sponzora} & Partneři
klubu \\
\textbf{Kontakty} & \texttt{kontakt} & \texttt{pozice},
\texttt{telefon}, \texttt{email}, \texttt{poradi} & Výbor klubu \\
\end{longtable}
}

Ukázka registrace CPT pro zápasy:

\begin{Shaded}
\begin{Highlighting}[]
\NormalTok{register\_post\_type(}\StringTok{\textquotesingle{}zapas\textquotesingle{}}\OtherTok{,} \DataTypeTok{array}\NormalTok{(}
    \StringTok{\textquotesingle{}labels\textquotesingle{}}\NormalTok{ =\textgreater{} }\DataTypeTok{array}\NormalTok{(}
        \StringTok{\textquotesingle{}name\textquotesingle{}}\NormalTok{               =\textgreater{} }\StringTok{\textquotesingle{}Zápasy\textquotesingle{}}\OtherTok{,}
        \StringTok{\textquotesingle{}singular\_name\textquotesingle{}}\NormalTok{      =\textgreater{} }\StringTok{\textquotesingle{}Zápas\textquotesingle{}}\OtherTok{,}
        \StringTok{\textquotesingle{}add\_new\textquotesingle{}}\NormalTok{            =\textgreater{} }\StringTok{\textquotesingle{}Přidat zápas\textquotesingle{}}\OtherTok{,}
        \StringTok{\textquotesingle{}add\_new\_item\textquotesingle{}}\NormalTok{       =\textgreater{} }\StringTok{\textquotesingle{}Přidat nový zápas\textquotesingle{}}\OtherTok{,}
        \CommentTok{// ...}
\NormalTok{    )}\OtherTok{,}
    \StringTok{\textquotesingle{}public\textquotesingle{}}\NormalTok{              =\textgreater{} }\KeywordTok{true}\OtherTok{,}
    \StringTok{\textquotesingle{}has\_archive\textquotesingle{}}\NormalTok{         =\textgreater{} }\KeywordTok{true}\OtherTok{,}
    \StringTok{\textquotesingle{}rewrite\textquotesingle{}}\NormalTok{             =\textgreater{} }\DataTypeTok{array}\NormalTok{(}\StringTok{\textquotesingle{}slug\textquotesingle{}}\NormalTok{ =\textgreater{} }\StringTok{\textquotesingle{}zapasy\textquotesingle{}}\NormalTok{)}\OtherTok{,}
    \StringTok{\textquotesingle{}supports\textquotesingle{}}\NormalTok{            =\textgreater{} }\DataTypeTok{array}\NormalTok{(}\StringTok{\textquotesingle{}title\textquotesingle{}}\OtherTok{,} \StringTok{\textquotesingle{}thumbnail\textquotesingle{}}\NormalTok{)}\OtherTok{,}
    \StringTok{\textquotesingle{}menu\_icon\textquotesingle{}}\NormalTok{           =\textgreater{} }\StringTok{\textquotesingle{}dashicons{-}awards\textquotesingle{}}\OtherTok{,}
    \StringTok{\textquotesingle{}show\_in\_rest\textquotesingle{}}\NormalTok{        =\textgreater{} }\KeywordTok{true}\OtherTok{,}
    \StringTok{\textquotesingle{}taxonomies\textquotesingle{}}\NormalTok{          =\textgreater{} }\DataTypeTok{array}\NormalTok{(}\StringTok{\textquotesingle{}sezona\textquotesingle{}}\OtherTok{,} \StringTok{\textquotesingle{}kategorie{-}tymu\textquotesingle{}}\OtherTok{,} \StringTok{\textquotesingle{}stav{-}zapasu\textquotesingle{}}\NormalTok{)}\OtherTok{,}
    \StringTok{\textquotesingle{}capability\_type\textquotesingle{}}\NormalTok{     =\textgreater{} }\DataTypeTok{array}\NormalTok{(}\StringTok{\textquotesingle{}zapas\textquotesingle{}}\OtherTok{,} \StringTok{\textquotesingle{}zapasy\textquotesingle{}}\NormalTok{)}\OtherTok{,}
    \StringTok{\textquotesingle{}map\_meta\_cap\textquotesingle{}}\NormalTok{        =\textgreater{} }\KeywordTok{true}\OtherTok{,}
\NormalTok{))}\OtherTok{;}
\end{Highlighting}
\end{Shaded}

Každý CPT má vlastní \texttt{capability\_type} s
\texttt{map\_meta\_cap\ =\textgreater{}\ true}, což umožňuje granulární
řízení oprávnění pro jednotlivé uživatelské role.

\paragraph{3.3.2 Taxonomie}\label{taxonomie}

Projekt používá čtyři vlastní taxonomie registrované přes
\texttt{register\_taxonomy()} {[}3{]}:

{\def\LTcaptype{none} % do not increment counter
\begin{longtable}[]{@{}
  >{\raggedright\arraybackslash}p{(\linewidth - 6\tabcolsep) * \real{0.2500}}
  >{\raggedright\arraybackslash}p{(\linewidth - 6\tabcolsep) * \real{0.2500}}
  >{\raggedright\arraybackslash}p{(\linewidth - 6\tabcolsep) * \real{0.2500}}
  >{\raggedright\arraybackslash}p{(\linewidth - 6\tabcolsep) * \real{0.2500}}@{}}
\toprule\noalign{}
\begin{minipage}[b]{\linewidth}\raggedright
Taxonomie
\end{minipage} & \begin{minipage}[b]{\linewidth}\raggedright
Slug
\end{minipage} & \begin{minipage}[b]{\linewidth}\raggedright
Sdílená mezi CPT
\end{minipage} & \begin{minipage}[b]{\linewidth}\raggedright
Účel
\end{minipage} \\
\midrule\noalign{}
\endhead
\bottomrule\noalign{}
\endlastfoot
\textbf{Sezóna} & \texttt{sezona} & zapas, tym, hrac, galerie &
Filtrování podle ročníku \\
\textbf{Kategorie týmu} & \texttt{kategorie-tymu} & zapas, tym, hrac,
galerie & Muži A, B, Dorost, Žáci\ldots{} \\
\textbf{Stav zápasu} & \texttt{stav-zapasu} & zapas & Nadcházející,
Odehraný, Zrušený \\
\textbf{Pozice hráče} & \texttt{pozice-hrace} & hrac & Brankáři, Hráči v
poli \\
\end{longtable}
}

Ukázka registrace taxonomie:

\begin{Shaded}
\begin{Highlighting}[]
\NormalTok{register\_taxonomy(}\StringTok{\textquotesingle{}sezona\textquotesingle{}}\OtherTok{,} \DataTypeTok{array}\NormalTok{(}\StringTok{\textquotesingle{}zapas\textquotesingle{}}\OtherTok{,} \StringTok{\textquotesingle{}tym\textquotesingle{}}\OtherTok{,} \StringTok{\textquotesingle{}hrac\textquotesingle{}}\OtherTok{,} \StringTok{\textquotesingle{}galerie\textquotesingle{}}\NormalTok{)}\OtherTok{,} \DataTypeTok{array}\NormalTok{(}
    \StringTok{\textquotesingle{}labels\textquotesingle{}}\NormalTok{            =\textgreater{} }\DataTypeTok{array}\NormalTok{(}
        \StringTok{\textquotesingle{}name\textquotesingle{}}\NormalTok{          =\textgreater{} }\StringTok{\textquotesingle{}Sezóny\textquotesingle{}}\OtherTok{,}
        \StringTok{\textquotesingle{}singular\_name\textquotesingle{}}\NormalTok{ =\textgreater{} }\StringTok{\textquotesingle{}Sezóna\textquotesingle{}}\OtherTok{,}
\NormalTok{    )}\OtherTok{,}
    \StringTok{\textquotesingle{}hierarchical\textquotesingle{}}\NormalTok{      =\textgreater{} }\KeywordTok{false}\OtherTok{,}
    \StringTok{\textquotesingle{}public\textquotesingle{}}\NormalTok{            =\textgreater{} }\KeywordTok{true}\OtherTok{,}
    \StringTok{\textquotesingle{}show\_admin\_column\textquotesingle{}}\NormalTok{ =\textgreater{} }\KeywordTok{true}\OtherTok{,}
    \StringTok{\textquotesingle{}show\_in\_rest\textquotesingle{}}\NormalTok{      =\textgreater{} }\KeywordTok{true}\OtherTok{,}
\NormalTok{))}\OtherTok{;}
\end{Highlighting}
\end{Shaded}

Sdílení taxonomií mezi více CPT umožňuje jednotnou klasifikaci obsahu a
zjednodušuje implementaci filtrování i návazných výpisů napříč
jednotlivými sekcemi webu.

\paragraph{3.3.3 Vazby mezi entitami}\label{vazby-mezi-entitami}

Vztahy mezi typy obsahu jsou řešeny dvěma mechanismy:

\textbf{Taxonomická vazba (M:N)} -- sezóna a kategorie týmu jsou sdíleny
mezi zápasy, týmy, hráči a galeriemi pomocí WordPress taxonomií.

\textbf{Meta pole vazba (1:N)} -- hráči jsou propojeni s týmy přes meta
pole \texttt{tym\_slug}. Každý hráč má uložen slug svého týmu (např.
\texttt{muzi-a}), který odpovídá hodnotě \texttt{tym\_slug} u
příslušného záznamu CPT \texttt{tym}.

\begin{verbatim}
Zápas ──(taxonomie)──> Sezóna <──(taxonomie)── Tým
  │                                               │
  └──(taxonomie)──> Kategorie týmu <──(taxonomie)─┘
                                                  │
                                            (meta: tym_slug)
                                                  │
                                                Hráč
\end{verbatim}

Funkce \texttt{slavoj\_count\_hracu\_tymu()} automaticky počítá hráče
přiřazené k týmu a výsledek ukládá do meta pole \texttt{pocet\_hracu}.

\paragraph{3.3.4 ER diagram}\label{er-diagram}

\begin{verbatim}
┌────────────────────────────────┐
│            SEZÓNA              │  taxonomie: sezona
│  slug: 2024-25 / 2025-26 …    │
└──────────┬─────────────────────┘
           │ m:n (term relationship)
     ┌─────┼──────┬─────────────────┐
     │     │      │                 │
     ▼     ▼      ▼                 ▼
┌────────┐ ┌────────┐ ┌──────────┐ ┌─────────┐
│ ZÁPAS  │ │  TÝM   │ │  HRÁČ    │ │ GALERIE │
└───┬────┘ └───┬────┘ └────┬─────┘ └─────────┘
    │          │           │
    │          │  1:n      │ n:1 (tym_slug)
    │          └───────────┘
    ▼ m:n
┌────────────────────────────────┐
│         KATEGORIE TÝMU         │  taxonomie: kategorie-tymu
└────────────────────────────────┘

┌─────────┐   ┌──────────┐
│ KONTAKT │   │ SPONZOR  │   (samostatné entity)
└─────────┘   └──────────┘
\end{verbatim}

\subsubsection{3.4 Vlastní meta pole (nahrazuje Advanced Custom
Fields)}\label{vlastnuxed-meta-pole-nahrazuje-advanced-custom-fields}

Místo pluginu Advanced Custom Fields {[}9{]} jsou meta boxy
implementovány přímo v PHP pomocí WordPress funkcí
\texttt{add\_meta\_box()}, \texttt{get\_post\_meta()} a
\texttt{update\_post\_meta()} {[}3{]}.

Ukázka definice meta boxu pro zápasy:

\begin{Shaded}
\begin{Highlighting}[]
\NormalTok{add\_meta\_box(}
    \StringTok{\textquotesingle{}slavoj\_zapas\_detail\textquotesingle{}}\OtherTok{,}
    \StringTok{\textquotesingle{}Detail zápasu\textquotesingle{}}\OtherTok{,}
    \StringTok{\textquotesingle{}slavoj\_zapas\_meta\_box\_html\textquotesingle{}}\OtherTok{,}
    \StringTok{\textquotesingle{}zapas\textquotesingle{}}\OtherTok{,}
    \StringTok{\textquotesingle{}normal\textquotesingle{}}\OtherTok{,}
    \StringTok{\textquotesingle{}high\textquotesingle{}}
\NormalTok{)}\OtherTok{;}
\end{Highlighting}
\end{Shaded}

Pole pro jednotlivé CPT:

\textbf{Zápas:} datum zápasu (date), čas výkopu (time), domácí tým
(text), hostující tým (text), skóre (text), střelci (text).

\textbf{Tým:} slug týmu (text), hlavní trenér (text), asistent trenéra
(text), zdravotník (text), počet hráčů (automaticky vypočteno).

\textbf{Hráč:} číslo dresu (number 1--99), rok narození (number), slug
týmu (text).

\textbf{Kontakt:} funkce/pozice (text), telefon (text), e-mail (email),
pořadí zobrazení (number).

\textbf{Sponzor:} webové stránky (url).

\subsubsection{3.5 Filtrace obsahu}\label{filtrace-obsahu}

Místo pluginů FacetWP nebo Search \& Filter Pro {[}10{]} jsou filtry
implementovány jako HTML formuláře s GET parametry:

\begin{Shaded}
\begin{Highlighting}[]
\DataTypeTok{\textless{}}\KeywordTok{form}\OtherTok{ method}\OperatorTok{=}\StringTok{"get"}\OtherTok{ action}\OperatorTok{=}\StringTok{"}\ErrorTok{\textless{}}\StringTok{?php echo esc\_url(get\_post\_type\_archive\_link(\textquotesingle{}zapas\textquotesingle{})); ?\textgreater{}"}\DataTypeTok{\textgreater{}}
    \DataTypeTok{\textless{}}\KeywordTok{select}\OtherTok{ name}\OperatorTok{=}\StringTok{"kat"}\OtherTok{ onchange}\OperatorTok{=}\StringTok{"this.form.submit()"}\DataTypeTok{\textgreater{}}
        \DataTypeTok{\textless{}}\KeywordTok{option}\OtherTok{ value}\OperatorTok{=}\StringTok{""}\DataTypeTok{\textgreater{}}\NormalTok{Všechny týmy}\DataTypeTok{\textless{}/}\KeywordTok{option}\DataTypeTok{\textgreater{}}
        \CommentTok{\textless{}!{-}{-} dynamicky generované možnosti {-}{-}\textgreater{}}
    \DataTypeTok{\textless{}/}\KeywordTok{select}\DataTypeTok{\textgreater{}}
    \DataTypeTok{\textless{}}\KeywordTok{noscript}\DataTypeTok{\textgreater{}\textless{}}\KeywordTok{button}\OtherTok{ type}\OperatorTok{=}\StringTok{"submit"}\DataTypeTok{\textgreater{}}\NormalTok{Filtrovat}\DataTypeTok{\textless{}/}\KeywordTok{button}\DataTypeTok{\textgreater{}\textless{}/}\KeywordTok{noscript}\DataTypeTok{\textgreater{}}
\DataTypeTok{\textless{}/}\KeywordTok{form}\DataTypeTok{\textgreater{}}
\end{Highlighting}
\end{Shaded}

PHP kód v šabloně načte GET parametry, sanitizuje je pomocí
\texttt{sanitize\_text\_field()} a předá do \texttt{WP\_Query} s
\texttt{tax\_query} a \texttt{meta\_query} {[}3{]}.

Šablony s filtrací:

{\def\LTcaptype{none} % do not increment counter
\begin{longtable}[]{@{}ll@{}}
\toprule\noalign{}
Šablona & Filtry \\
\midrule\noalign{}
\endhead
\bottomrule\noalign{}
\endlastfoot
\texttt{archive-zapas.php} & Tým, sezóna, stav zápasu \\
\texttt{archive-tym.php} & Sezóna \\
\texttt{archive-galerie.php} & Tým, sezóna \\
\end{longtable}
}

Formuláře používají \texttt{onchange="this.form.submit()"} pro
automatické odeslání při změně výběru, s
\texttt{\textless{}noscript\textgreater{}} fallbackem pro prohlížeče bez
JavaScriptu.

\subsubsection{3.6 Řazení a kontextové filtrování
taxonomií}\label{ux159azenuxed-a-kontextovuxe9-filtrovuxe1nuxed-taxonomiuxed}

\paragraph{3.6.1 Problém}\label{probluxe9m}

WordPress řadí termy taxonomií abecedně. Pro fotbalový klub je přirozené
hierarchické pořadí: Muži A → Muži B → Dorost → Starší žáci → Mladší
žáci → Přípravky. Navíc kategorie „Stará garda'' se má zobrazovat pouze
u galerií, nikoli u zápasů nebo hráčů.

\paragraph{3.6.2 Řešení}\label{ux159eux161enuxed}

Pořadí je definováno centrálně ve funkci
\texttt{slavoj\_kategorie\_poradi()}:

\begin{Shaded}
\begin{Highlighting}[]
\KeywordTok{function}\NormalTok{ slavoj\_kategorie\_poradi() \{}
    \ControlFlowTok{return} \DataTypeTok{array}\NormalTok{(}
        \StringTok{\textquotesingle{}muzi{-}a\textquotesingle{}}\NormalTok{            =\textgreater{} }\StringTok{\textquotesingle{}Muži A\textquotesingle{}}\OtherTok{,}
        \StringTok{\textquotesingle{}muzi{-}b\textquotesingle{}}\NormalTok{            =\textgreater{} }\StringTok{\textquotesingle{}Muži B\textquotesingle{}}\OtherTok{,}
        \StringTok{\textquotesingle{}stara{-}garda\textquotesingle{}}\NormalTok{       =\textgreater{} }\StringTok{\textquotesingle{}Stará garda\textquotesingle{}}\OtherTok{,}
        \StringTok{\textquotesingle{}dorost\textquotesingle{}}\NormalTok{            =\textgreater{} }\StringTok{\textquotesingle{}Dorost\textquotesingle{}}\OtherTok{,}
        \StringTok{\textquotesingle{}starsi{-}zaci\textquotesingle{}}\NormalTok{       =\textgreater{} }\StringTok{\textquotesingle{}Starší žáci\textquotesingle{}}\OtherTok{,}
        \StringTok{\textquotesingle{}mladsi{-}zaci\textquotesingle{}}\NormalTok{       =\textgreater{} }\StringTok{\textquotesingle{}Mladší žáci\textquotesingle{}}\OtherTok{,}
        \StringTok{\textquotesingle{}starsi{-}pripravka\textquotesingle{}}\NormalTok{  =\textgreater{} }\StringTok{\textquotesingle{}Starší přípravka\textquotesingle{}}\OtherTok{,}
        \StringTok{\textquotesingle{}mladsi{-}pripravka\textquotesingle{}}\NormalTok{  =\textgreater{} }\StringTok{\textquotesingle{}Mladší přípravka\textquotesingle{}}\OtherTok{,}
        \StringTok{\textquotesingle{}mini{-}pripravka\textquotesingle{}}\NormalTok{    =\textgreater{} }\StringTok{\textquotesingle{}Mini přípravka\textquotesingle{}}\OtherTok{,}
\NormalTok{    )}\OtherTok{;}
\NormalTok{\}}
\end{Highlighting}
\end{Shaded}

Globální filtr na WordPress hook \texttt{get\_terms} automaticky seřadí
termy podle tohoto kanonického pořadí a odfiltruje kontextově omezené
termy (např. „Stará garda'' se zobrazí pouze v kontextu galerií).

\subsubsection{3.7 Šablony}\label{ux161ablony}

\paragraph{3.7.1 Přehled šablon}\label{pux159ehled-ux161ablon}

{\def\LTcaptype{none} % do not increment counter
\begin{longtable}[]{@{}lll@{}}
\toprule\noalign{}
Šablona & URL & Obsah \\
\midrule\noalign{}
\endhead
\bottomrule\noalign{}
\endlastfoot
\texttt{front-page.php} & \texttt{/} & Nadcházející zápasy, aktuality \\
\texttt{archive-zapas.php} & \texttt{/zapasy/} & Filtrovaný seznam
zápasů \\
\texttt{archive-tym.php} & \texttt{/tymy/} & Karty týmů \\
\texttt{archive-galerie.php} & \texttt{/galerie/} & Grid fotoalb \\
\texttt{single-zapas.php} & \texttt{/zapasy/\{slug\}/} & Detail
zápasu \\
\texttt{single-tym.php} & \texttt{/tymy/\{slug\}/} & Profil týmu,
soupiska \\
\texttt{single-hrac.php} & \texttt{/hraci/\{slug\}/} & Profil hráče \\
\texttt{single-galerie.php} & \texttt{/galerie/\{slug\}/} & Fotoalbum s
lightboxem \\
\texttt{page-historie.php} & \texttt{/historie/} & Historie klubu \\
\texttt{page-kontakty.php} & \texttt{/kontakty/} & Kontakty výboru \\
\texttt{page-sponzori.php} & \texttt{/sponzori/} & Loga partnerů \\
\end{longtable}
}

\paragraph{3.7.2 Znovupoužitelné komponenty
(template-parts)}\label{znovupouux17eitelnuxe9-komponenty-template-parts}

\begin{itemize}
\tightlist
\item
  \textbf{\texttt{card-match.php}} -- karta zápasu (domácí vs hosté,
  datum, skóre, střelci)
\item
  \textbf{\texttt{hero-team.php}} -- hero sekce s názvem týmu a trenéry
\item
  \textbf{\texttt{site-header.php}} -- hlavička webu s navigací
\item
  \textbf{\texttt{site-footer.php}} -- patička webu
\end{itemize}

\subsubsection{3.8 CSS architektura}\label{css-architektura}

Styly jsou rozdělené do komponentových souborů a načítány přes
\texttt{wp\_enqueue\_style()} {[}3{]} s definovanými závislostmi:

\begin{Shaded}
\begin{Highlighting}[]
\KeywordTok{function}\NormalTok{ slavoj\_enqueue\_scripts() \{}
    \VariableTok{$ver} \OperatorTok{=}\NormalTok{ wp\_get\_theme(){-}\textgreater{}get(}\StringTok{\textquotesingle{}Version\textquotesingle{}}\NormalTok{)}\OtherTok{;}
\NormalTok{    wp\_enqueue\_style(}\StringTok{\textquotesingle{}bootstrap\textquotesingle{}}\OtherTok{,} \StringTok{\textquotesingle{}https://cdn.jsdelivr.net/npm/bootstrap@5.3.3/dist/css/bootstrap.min.css\textquotesingle{}}\OtherTok{,} \DataTypeTok{array}\NormalTok{()}\OtherTok{,} \StringTok{\textquotesingle{}5.3.3\textquotesingle{}}\NormalTok{)}\OtherTok{;}
\NormalTok{    wp\_enqueue\_script(}\StringTok{\textquotesingle{}bootstrap{-}js\textquotesingle{}}\OtherTok{,} \StringTok{\textquotesingle{}https://cdn.jsdelivr.net/npm/bootstrap@5.3.3/dist/js/bootstrap.bundle.min.js\textquotesingle{}}\OtherTok{,} \DataTypeTok{array}\NormalTok{()}\OtherTok{,} \StringTok{\textquotesingle{}5.3.3\textquotesingle{}}\OtherTok{,} \KeywordTok{true}\NormalTok{)}\OtherTok{;}

    \VariableTok{$css} \OperatorTok{=}\NormalTok{ get\_template\_directory\_uri() }\OperatorTok{.} \StringTok{\textquotesingle{}/assets/css/\textquotesingle{}}\OtherTok{;}
\NormalTok{    wp\_enqueue\_style(}\StringTok{\textquotesingle{}slavoj{-}utilities\textquotesingle{}}\OtherTok{,}  \VariableTok{$css} \OperatorTok{.} \StringTok{\textquotesingle{}utilities.css\textquotesingle{}}\OtherTok{,}            \DataTypeTok{array}\NormalTok{(}\StringTok{\textquotesingle{}bootstrap\textquotesingle{}}\NormalTok{)}\OtherTok{,} \VariableTok{$ver}\NormalTok{)}\OtherTok{;}
\NormalTok{    wp\_enqueue\_style(}\StringTok{\textquotesingle{}slavoj{-}header\textquotesingle{}}\OtherTok{,}     \VariableTok{$css} \OperatorTok{.} \StringTok{\textquotesingle{}components/header.css\textquotesingle{}}\OtherTok{,}    \DataTypeTok{array}\NormalTok{(}\StringTok{\textquotesingle{}slavoj{-}utilities\textquotesingle{}}\NormalTok{)}\OtherTok{,} \VariableTok{$ver}\NormalTok{)}\OtherTok{;}
\NormalTok{    wp\_enqueue\_style(}\StringTok{\textquotesingle{}slavoj{-}buttons\textquotesingle{}}\OtherTok{,}    \VariableTok{$css} \OperatorTok{.} \StringTok{\textquotesingle{}components/buttons.css\textquotesingle{}}\OtherTok{,}   \DataTypeTok{array}\NormalTok{(}\StringTok{\textquotesingle{}slavoj{-}utilities\textquotesingle{}}\NormalTok{)}\OtherTok{,} \VariableTok{$ver}\NormalTok{)}\OtherTok{;}
\NormalTok{    wp\_enqueue\_style(}\StringTok{\textquotesingle{}slavoj{-}cards\textquotesingle{}}\OtherTok{,}      \VariableTok{$css} \OperatorTok{.} \StringTok{\textquotesingle{}components/cards.css\textquotesingle{}}\OtherTok{,}     \DataTypeTok{array}\NormalTok{(}\StringTok{\textquotesingle{}slavoj{-}utilities\textquotesingle{}}\NormalTok{)}\OtherTok{,} \VariableTok{$ver}\NormalTok{)}\OtherTok{;}
\NormalTok{    wp\_enqueue\_style(}\StringTok{\textquotesingle{}slavoj{-}main\textquotesingle{}}\OtherTok{,} \VariableTok{$css} \OperatorTok{.} \StringTok{\textquotesingle{}main.css\textquotesingle{}}\OtherTok{,} \DataTypeTok{array}\NormalTok{(}\StringTok{\textquotesingle{}slavoj{-}utilities\textquotesingle{}}\OtherTok{,} \StringTok{\textquotesingle{}slavoj{-}buttons\textquotesingle{}}\OtherTok{,} \StringTok{\textquotesingle{}slavoj{-}cards\textquotesingle{}}\OtherTok{,} \StringTok{\textquotesingle{}slavoj{-}header\textquotesingle{}}\NormalTok{)}\OtherTok{,} \VariableTok{$ver}\NormalTok{)}\OtherTok{;}
\NormalTok{\}}
\NormalTok{add\_action(}\StringTok{\textquotesingle{}wp\_enqueue\_scripts\textquotesingle{}}\OtherTok{,} \StringTok{\textquotesingle{}slavoj\_enqueue\_scripts\textquotesingle{}}\NormalTok{)}\OtherTok{;}
\end{Highlighting}
\end{Shaded}

Soubor \texttt{style.css} v kořeni tématu slouží pouze jako
identifikátor pro WordPress (název tématu, verze, autor) a neobsahuje
žádné styly. Veškeré styly a skripty jsou načítány přes
\texttt{wp\_enqueue\_style()} a \texttt{wp\_enqueue\_script()} -- nikdy
přímo v šablonách {[}3{]}.

\begin{verbatim}
assets/css/
├── utilities.css              # Barvy, spacing, utility třídy
├── main.css                   # Hlavní styly šablon
└── components/
    ├── header.css             # Navigace a hlavička
    ├── nav-mobile.css         # Mobilní menu
    ├── buttons.css            # Tlačítka
    ├── hero.css               # Hero sekce
    └── cards.css              # Karty zápasů, galerií, hráčů
\end{verbatim}

\subsubsection{3.9 Plugin Slavoj Custom
Fields}\label{plugin-slavoj-custom-fields}

Plugin je uložen v souboru \texttt{slavoj-custom-fields.php} a poskytuje
administrační nástroje:

\textbf{Vlastní sloupce v přehledech:} - Zápasy: datum (řaditelný
chronologicky), tým, skóre - Hráči: číslo dresu (řaditelné numericky),
pozice, tým

Ukázka implementace vlastního sloupce:

\begin{Shaded}
\begin{Highlighting}[]
\KeywordTok{function}\NormalTok{ slavoj\_zapas\_admin\_column\_data(}\VariableTok{$column}\OtherTok{,} \VariableTok{$post\_id}\NormalTok{) \{}
    \ControlFlowTok{if}\NormalTok{ (}\VariableTok{$column} \OperatorTok{===} \StringTok{\textquotesingle{}datum\_zapasu\textquotesingle{}}\NormalTok{) \{}
        \VariableTok{$datum} \OperatorTok{=}\NormalTok{ get\_post\_meta(}\VariableTok{$post\_id}\OtherTok{,} \StringTok{\textquotesingle{}datum\_zapasu\textquotesingle{}}\OtherTok{,} \KeywordTok{true}\NormalTok{)}\OtherTok{;}
        \ControlFlowTok{if}\NormalTok{ (}\VariableTok{$datum}\NormalTok{) \{}
            \VariableTok{$dt} \OperatorTok{=} \BuiltInTok{DateTime}\NormalTok{::createFromFormat(}\StringTok{\textquotesingle{}Y{-}m{-}d\textquotesingle{}}\OtherTok{,} \VariableTok{$datum}\NormalTok{)}\OtherTok{;}
            \KeywordTok{echo} \VariableTok{$dt} \OtherTok{?}\NormalTok{ esc\_html(}\VariableTok{$dt}\NormalTok{{-}\textgreater{}format(}\StringTok{\textquotesingle{}j. n. Y\textquotesingle{}}\NormalTok{)) }\OtherTok{:}\NormalTok{ esc\_html(}\VariableTok{$datum}\NormalTok{)}\OtherTok{;}
\NormalTok{        \}}
\NormalTok{    \}}
    \CommentTok{// ...}
\NormalTok{\}}
\end{Highlighting}
\end{Shaded}

\textbf{Dropdown filtry:} nad seznamy příspěvků lze filtrovat podle
sezóny, kategorie týmu nebo týmu (u hráčů).

\textbf{Nástrojová stránka} (Nástroje → Slavoj nastavení): - přehled
registrovaných CPT a taxonomií - tlačítko „Vytvořit výchozí hodnoty
taxonomií'' - tlačítko „Vytvořit ukázková data'' (zápasy, hráči,
kontakty, galerie) - tlačítko „Vytvořit stránky webu''

\subsubsection{3.10 Uživatelské role}\label{uux17eivatelskuxe9-role}

Místo pluginu User Role Editor {[}11{]} je vlastní role registrována
přímo v kódu:

\begin{Shaded}
\begin{Highlighting}[]
\NormalTok{add\_role(}\StringTok{\textquotesingle{}slavoj\_editor\textquotesingle{}}\OtherTok{,} \StringTok{\textquotesingle{}Správce obsahu\textquotesingle{}}\OtherTok{,} \DataTypeTok{array}\NormalTok{(}
    \StringTok{\textquotesingle{}read\textquotesingle{}}\NormalTok{                   =\textgreater{} }\KeywordTok{true}\OtherTok{,}
    \StringTok{\textquotesingle{}edit\_posts\textquotesingle{}}\NormalTok{             =\textgreater{} }\KeywordTok{true}\OtherTok{,}
    \StringTok{\textquotesingle{}edit\_others\_posts\textquotesingle{}}\NormalTok{      =\textgreater{} }\KeywordTok{true}\OtherTok{,}
    \StringTok{\textquotesingle{}publish\_posts\textquotesingle{}}\NormalTok{          =\textgreater{} }\KeywordTok{true}\OtherTok{,}
    \StringTok{\textquotesingle{}upload\_files\textquotesingle{}}\NormalTok{           =\textgreater{} }\KeywordTok{true}\OtherTok{,}
    \StringTok{\textquotesingle{}manage\_categories\textquotesingle{}}\NormalTok{      =\textgreater{} }\KeywordTok{true}\OtherTok{,}
    \CommentTok{// + oprávnění pro všechny CPT}
\NormalTok{))}\OtherTok{;}
\end{Highlighting}
\end{Shaded}

{\def\LTcaptype{none} % do not increment counter
\begin{longtable}[]{@{}
  >{\raggedright\arraybackslash}p{(\linewidth - 4\tabcolsep) * \real{0.2500}}
  >{\raggedright\arraybackslash}p{(\linewidth - 4\tabcolsep) * \real{0.2917}}
  >{\raggedright\arraybackslash}p{(\linewidth - 4\tabcolsep) * \real{0.4583}}@{}}
\toprule\noalign{}
\begin{minipage}[b]{\linewidth}\raggedright
Role
\end{minipage} & \begin{minipage}[b]{\linewidth}\raggedright
Popis
\end{minipage} & \begin{minipage}[b]{\linewidth}\raggedright
Oprávnění
\end{minipage} \\
\midrule\noalign{}
\endhead
\bottomrule\noalign{}
\endlastfoot
Administrátor & Správce webu & Plný přístup \\
Správce obsahu & Správce zápasů a obsahu & Editace CPT, nahrávání médií,
bez nastavení WP \\
\end{longtable}
}

\subsubsection{3.11 Optimalizace výkonu}\label{optimalizace-vuxfdkonu}

Základní optimalizace jsou implementovány přímo v kódu tématu:

{\def\LTcaptype{none} % do not increment counter
\begin{longtable}[]{@{}
  >{\raggedright\arraybackslash}p{(\linewidth - 2\tabcolsep) * \real{0.5000}}
  >{\raggedright\arraybackslash}p{(\linewidth - 2\tabcolsep) * \real{0.5000}}@{}}
\toprule\noalign{}
\begin{minipage}[b]{\linewidth}\raggedright
Optimalizace
\end{minipage} & \begin{minipage}[b]{\linewidth}\raggedright
Implementace
\end{minipage} \\
\midrule\noalign{}
\endhead
\bottomrule\noalign{}
\endlastfoot
Lazy loading obrázků & Funkce \texttt{slavoj\_add\_lazy\_loading()}
přidává \texttt{loading="lazy"} {[}12{]} \\
Bootstrap z CDN & CSS a JS z jsdelivr.net -- využívá cache prohlížeče
{[}5{]} \\
Minimální JavaScript & Filtry přes HTML formuláře, bez AJAX knihoven \\
Mobile-first CSS & Vlastní CSS psaný mobile-first, minimální velikost \\
\end{longtable}
}

\begin{Shaded}
\begin{Highlighting}[]
\KeywordTok{function}\NormalTok{ slavoj\_add\_lazy\_loading(}\VariableTok{$content}\NormalTok{) \{}
    \ControlFlowTok{return} \FunctionTok{str\_replace}\NormalTok{(}\StringTok{\textquotesingle{}\textless{}img \textquotesingle{}}\OtherTok{,} \StringTok{\textquotesingle{}\textless{}img loading="lazy" \textquotesingle{}}\OtherTok{,} \VariableTok{$content}\NormalTok{)}\OtherTok{;}
\NormalTok{\}}
\NormalTok{add\_filter(}\StringTok{\textquotesingle{}the\_content\textquotesingle{}}\OtherTok{,} \StringTok{\textquotesingle{}slavoj\_add\_lazy\_loading\textquotesingle{}}\NormalTok{)}\OtherTok{;}
\NormalTok{add\_filter(}\StringTok{\textquotesingle{}post\_thumbnail\_html\textquotesingle{}}\OtherTok{,} \StringTok{\textquotesingle{}slavoj\_add\_lazy\_loading\textquotesingle{}}\NormalTok{)}\OtherTok{;}
\end{Highlighting}
\end{Shaded}

\subsubsection{3.12 Bezpečnost}\label{bezpeux10dnost}

Projekt dodržuje standardní bezpečnostní praktiky WordPressu {[}3{]}
{[}13{]}:

\begin{itemize}
\tightlist
\item
  \textbf{Sanitizace vstupů:} Všechna data z formulářů procházejí přes
  \texttt{sanitize\_text\_field()} nebo \texttt{sanitize\_email()}.
\item
  \textbf{Escapování výstupů:} Veškerý výstup do HTML používá
  \texttt{esc\_html()}, \texttt{esc\_attr()} nebo \texttt{esc\_url()}.
\item
  \textbf{Nonce ochrana:} Každý meta box formulář je chráněn WordPress
  nonce tokenem (\texttt{wp\_nonce\_field} /
  \texttt{wp\_verify\_nonce}).
\item
  \textbf{WP\_Query:} Databázové dotazy využívají výhradně WordPress
  API, nikoli přímé SQL.
\item
  \textbf{Capability mapping:} Každý CPT má vlastní
  \texttt{capability\_type} pro granulární řízení oprávnění.
\item
  \textbf{WPS Hide Login:} Bezplatný plugin z oficiálního repozitáře
  WordPressu, který mění výchozí přihlašovací URL \texttt{/wp-login.php}
  na vlastní adresu (např. \texttt{/prihlaseni/}). Přístup na
  \texttt{/wp-login.php} i \texttt{/wp-admin/} (pro nepřihlášené) vrací
  chybu 404. Tím se výrazně snižuje riziko automatizovaných brute-force
  útoků, protože útočníci nemohou najít standardní přihlašovací stránku.
  Nastavení pluginu: \textbf{Nastavení → WPS Hide Login} -- stačí zadat
  požadovaný slug a uložit.
\end{itemize}

\subsubsection{3.13 JavaScript}\label{javascript}

Projekt záměrně minimalizuje použití JavaScriptu. Celkový rozsah
vlastního JS je přibližně 70 řádků {[}14{]}:

\begin{enumerate}
\def\labelenumi{\arabic{enumi}.}
\tightlist
\item
  \textbf{Filtry} -- \texttt{onchange="this.form.submit()"} s
  \texttt{\textless{}noscript\textgreater{}} fallbackem (1 řádek na
  formulář)
\item
  \textbf{Lightbox galerie} v \texttt{single-galerie.php}
  (\textasciitilde55 řádků) -- kliknutí na foto, navigace šipkami/Escape
\item
  \textbf{Quick-add střelců} v \texttt{single-zapas.php}
  (\textasciitilde12 řádků) -- administrační pomůcka
\end{enumerate}

Všechny funkce mají funkční fallback bez JavaScriptu. Žádné frameworky
ani build nástroje nejsou použity.

\begin{center}\rule{0.5\linewidth}{0.5pt}\end{center}

\subsection{4 Testování}\label{testovuxe1nuxed}

\subsubsection{4.1 Přehled testování}\label{pux159ehled-testovuxe1nuxed}

Testování probíhalo průběžně během vývoje. Byly provedeny následující
typy testů:

\begin{enumerate}
\def\labelenumi{\arabic{enumi}.}
\tightlist
\item
  \textbf{Funkční testování} -- ověření správné funkčnosti všech
  stránek, filtrů a administrace
\item
  \textbf{Bezpečnostní kontrola kódu} -- kontrola sanitizace, escapování
  a nonce ochrany
\item
  \textbf{Responzivní testování} -- ověření zobrazení na různých
  zařízeních
\item
  \textbf{Kontrola WordPress best practices} -- dodržení standardů pro
  vývoj témat {[}13{]}
\end{enumerate}

\subsubsection{4.2 Funkční testovací
scénáře}\label{funkux10dnuxed-testovacuxed-scuxe9nuxe1ux159e}

\paragraph{Scénář 1: Přidání nového
zápasu}\label{scuxe9nuxe1ux159-1-pux159iduxe1nuxed-novuxe9ho-zuxe1pasu}

{\def\LTcaptype{none} % do not increment counter
\begin{longtable}[]{@{}
  >{\raggedright\arraybackslash}p{(\linewidth - 6\tabcolsep) * \real{0.1579}}
  >{\raggedright\arraybackslash}p{(\linewidth - 6\tabcolsep) * \real{0.1579}}
  >{\raggedright\arraybackslash}p{(\linewidth - 6\tabcolsep) * \real{0.5263}}
  >{\raggedright\arraybackslash}p{(\linewidth - 6\tabcolsep) * \real{0.1579}}@{}}
\toprule\noalign{}
\begin{minipage}[b]{\linewidth}\raggedright
Krok
\end{minipage} & \begin{minipage}[b]{\linewidth}\raggedright
Akce
\end{minipage} & \begin{minipage}[b]{\linewidth}\raggedright
Očekávaný výsledek
\end{minipage} & \begin{minipage}[b]{\linewidth}\raggedright
Stav
\end{minipage} \\
\midrule\noalign{}
\endhead
\bottomrule\noalign{}
\endlastfoot
1 & Přihlášení do administrace & Zobrazí se dashboard & Ano \\
2 & Klik na Zápasy → Přidat zápas & Zobrazí se formulář & Ano \\
3 & Vyplnění polí (datum, čas, týmy, skóre) & Pole přijímají data &
Ano \\
4 & Přiřazení sezóny a kategorie týmu & Taxonomie přiřazeny & Ano \\
5 & Klik na Publikovat & Zápas uložen & Ano \\
6 & Zobrazení na \texttt{/zapasy/} & Zápas se zobrazí v seznamu & Ano \\
7 & Klik na detail zápasu & Detail s daty se zobrazí & Ano \\
\end{longtable}
}

\paragraph{Scénář 2: Filtrování
zápasů}\label{scuxe9nuxe1ux159-2-filtrovuxe1nuxed-zuxe1pasux16f}

{\def\LTcaptype{none} % do not increment counter
\begin{longtable}[]{@{}
  >{\raggedright\arraybackslash}p{(\linewidth - 6\tabcolsep) * \real{0.1579}}
  >{\raggedright\arraybackslash}p{(\linewidth - 6\tabcolsep) * \real{0.1579}}
  >{\raggedright\arraybackslash}p{(\linewidth - 6\tabcolsep) * \real{0.5263}}
  >{\raggedright\arraybackslash}p{(\linewidth - 6\tabcolsep) * \real{0.1579}}@{}}
\toprule\noalign{}
\begin{minipage}[b]{\linewidth}\raggedright
Krok
\end{minipage} & \begin{minipage}[b]{\linewidth}\raggedright
Akce
\end{minipage} & \begin{minipage}[b]{\linewidth}\raggedright
Očekávaný výsledek
\end{minipage} & \begin{minipage}[b]{\linewidth}\raggedright
Stav
\end{minipage} \\
\midrule\noalign{}
\endhead
\bottomrule\noalign{}
\endlastfoot
1 & Otevření \texttt{/zapasy/} & Zobrazí se všechny zápasy & Ano \\
2 & Výběr „Muži A'' ve filtru týmu & Stránka se přenačte, zobrazí pouze
Muži A & Ano \\
3 & Výběr „2025/26'' ve filtru sezóny & Zobrazí pouze zápasy dané sezóny
& Ano \\
4 & Výběr „Odehrané'' ve filtru stavu & Zobrazí pouze odehrané zápasy &
Ano \\
5 & Reset filtrů (výběr „Všechny'') & Zobrazí se všechny zápasy & Ano \\
\end{longtable}
}

\paragraph{Scénář 3: Soupiska
týmu}\label{scuxe9nuxe1ux159-3-soupiska-tuxfdmu}

{\def\LTcaptype{none} % do not increment counter
\begin{longtable}[]{@{}
  >{\raggedright\arraybackslash}p{(\linewidth - 6\tabcolsep) * \real{0.1579}}
  >{\raggedright\arraybackslash}p{(\linewidth - 6\tabcolsep) * \real{0.1579}}
  >{\raggedright\arraybackslash}p{(\linewidth - 6\tabcolsep) * \real{0.5263}}
  >{\raggedright\arraybackslash}p{(\linewidth - 6\tabcolsep) * \real{0.1579}}@{}}
\toprule\noalign{}
\begin{minipage}[b]{\linewidth}\raggedright
Krok
\end{minipage} & \begin{minipage}[b]{\linewidth}\raggedright
Akce
\end{minipage} & \begin{minipage}[b]{\linewidth}\raggedright
Očekávaný výsledek
\end{minipage} & \begin{minipage}[b]{\linewidth}\raggedright
Stav
\end{minipage} \\
\midrule\noalign{}
\endhead
\bottomrule\noalign{}
\endlastfoot
1 & Otevření \texttt{/tymy/} & Zobrazí se karty týmů & Ano \\
2 & Klik na tým & Detail týmu se soupiskou & Ano \\
3 & Ověření počtu hráčů & Odpovídá počtu hráčů v databázi & Ano \\
4 & Kontrola řazení hráčů & Seřazeni podle čísla dresu vzestupně &
Ano \\
\end{longtable}
}

\paragraph{Scénář 4: Galerie}\label{scuxe9nuxe1ux159-4-galerie}

{\def\LTcaptype{none} % do not increment counter
\begin{longtable}[]{@{}llll@{}}
\toprule\noalign{}
Krok & Akce & Očekávaný výsledek & Stav \\
\midrule\noalign{}
\endhead
\bottomrule\noalign{}
\endlastfoot
1 & Otevření \texttt{/galerie/} & Zobrazí se přehled alb & Ano \\
2 & Klik na album & Detail s fotogalerií & Ano \\
3 & Klik na fotografii & Lightbox se otevře & Ano \\
4 & Navigace šipkami/Escape & Funkční navigace, zavření & Ano \\
\end{longtable}
}

\paragraph{Scénář 5: Uživatelská role -- Správce
obsahu}\label{scuxe9nuxe1ux159-5-uux17eivatelskuxe1-role-spruxe1vce-obsahu}

{\def\LTcaptype{none} % do not increment counter
\begin{longtable}[]{@{}
  >{\raggedright\arraybackslash}p{(\linewidth - 6\tabcolsep) * \real{0.1579}}
  >{\raggedright\arraybackslash}p{(\linewidth - 6\tabcolsep) * \real{0.1579}}
  >{\raggedright\arraybackslash}p{(\linewidth - 6\tabcolsep) * \real{0.5263}}
  >{\raggedright\arraybackslash}p{(\linewidth - 6\tabcolsep) * \real{0.1579}}@{}}
\toprule\noalign{}
\begin{minipage}[b]{\linewidth}\raggedright
Krok
\end{minipage} & \begin{minipage}[b]{\linewidth}\raggedright
Akce
\end{minipage} & \begin{minipage}[b]{\linewidth}\raggedright
Očekávaný výsledek
\end{minipage} & \begin{minipage}[b]{\linewidth}\raggedright
Stav
\end{minipage} \\
\midrule\noalign{}
\endhead
\bottomrule\noalign{}
\endlastfoot
1 & Přihlášení jako Správce obsahu & Dashboard bez nastavení WP & Ano \\
2 & Editace zápasu & Povoleno & Ano \\
3 & Nahrání obrázku & Povoleno & Ano \\
4 & Přístup k Nastavení → Obecné & Zamítnuto & Ano \\
5 & Přístup k Pluginy & Zamítnuto & Ano \\
\end{longtable}
}

\subsubsection{4.3 Bezpečnostní
kontrola}\label{bezpeux10dnostnuxed-kontrola}

Kontrola kódu provedená v březnu 2026 potvrdila dodržení bezpečnostních
standardů {[}13{]}:

\begin{itemize}
\tightlist
\item
  Ano Sanitizace vstupů: \texttt{sanitize\_text\_field()},
  \texttt{sanitize\_email()}, \texttt{wp\_unslash()}
\item
  Ano Escapování výstupů: \texttt{esc\_html()}, \texttt{esc\_attr()},
  \texttt{esc\_url()}
\item
  Ano Nonce ochrana formulářů: \texttt{wp\_nonce\_field()},
  \texttt{wp\_verify\_nonce()}
\item
  Ano WP\_Query místo raw SQL
\item
  Ano Správné použití \texttt{wp\_reset\_postdata()} po custom queries
\end{itemize}

\subsubsection{4.4 Nalezené a opravené
chyby}\label{nalezenuxe9-a-opravenuxe9-chyby}

Během testování byly nalezeny a opraveny tyto chyby:

\begin{enumerate}
\def\labelenumi{\arabic{enumi}.}
\tightlist
\item
  \textbf{\texttt{single-galerie.php}} -- sezóna se načítala z meta pole
  místo z taxonomie. Opraveno na čtení z taxonomie \texttt{sezona}.
\item
  \textbf{\texttt{front-page.php}} -- Bootstrap třída \texttt{.card}
  přidávala nechtěný rámeček. Opraveno přidáním CSS
  \texttt{border:\ none}.
\item
  \textbf{\texttt{front-page.php}} -- horizontální scrollbar na
  mobilních zařízeních. Opraveno přidáním Bootstrap třídy \texttt{g-0}
  pro odstranění gutterů.
\end{enumerate}

\begin{center}\rule{0.5\linewidth}{0.5pt}\end{center}

\subsection{5 Příručka
administrátora}\label{pux159uxedruux10dka-administruxe1tora}

\subsubsection{5.1 Návod pro správce -- administrace
WordPress}\label{nuxe1vod-pro-spruxe1vce-administrace-wordpress}

\paragraph{5.1.1 Přihlášení}\label{pux159ihluxe1ux161enuxed}

Otevřete přihlašovací stránku webu a přihlaste se. Výchozí adresa
WordPressu je \texttt{/wp-login.php}, ale na produkčním serveru je
pomocí pluginu \textbf{WPS Hide Login} změněna na vlastní slug (viz
sekce 3.12). V levém menu uvidíte sekce: \textbf{Zápasy}, \textbf{Týmy},
\textbf{Hráči}, \textbf{Galerie}, \textbf{Sponzoři}, \textbf{Kontakty}.

\paragraph{5.1.2 Správa zápasů}\label{spruxe1va-zuxe1pasux16f}

\begin{enumerate}
\def\labelenumi{\arabic{enumi}.}
\tightlist
\item
  Klikněte \textbf{Zápasy → Přidat zápas}
\item
  Vyplňte titulek (např. „TJ Slavoj Mýto vs SK Nepomuk'')
\item
  V sekci \textbf{Detail zápasu} vyplňte datum, čas, domácí tým, hosty,
  skóre a střelce
\item
  V pravém panelu přiřaďte \textbf{kategorii týmu} a \textbf{sezónu}
\item
  Klikněte \textbf{Publikovat}
\end{enumerate}

Pro zadání výsledku odehraného zápasu: otevřete existující zápas,
vyplňte \textbf{Skóre} a \textbf{Střelce}, klikněte
\textbf{Aktualizovat}.

\paragraph{5.1.3 Správa týmů}\label{spruxe1va-tuxfdmux16f}

\begin{enumerate}
\def\labelenumi{\arabic{enumi}.}
\tightlist
\item
  Klikněte \textbf{Týmy → Přidat tým}
\item
  Vyplňte název, identifikátor týmu (slug, např. \texttt{muzi-a}),
  trenéry
\item
  Nastavte \textbf{Obrázek příspěvku} (týmová fotografie)
\item
  Přiřaďte sezónu a kategorii
\item
  \textbf{Počet hráčů} se vypočítá automaticky z přiřazených hráčů
\end{enumerate}

\paragraph{5.1.4 Správa hráčů}\label{spruxe1va-hruxe1ux10dux16f}

\begin{enumerate}
\def\labelenumi{\arabic{enumi}.}
\tightlist
\item
  Klikněte \textbf{Hráči → Přidat hráče}
\item
  Vyplňte jméno, číslo dresu (1--99), rok narození, slug týmu
\item
  Přiřaďte pozici (Brankář / Hráč v poli) a sezónu
\end{enumerate}

\paragraph{5.1.5 Správa galerií}\label{spruxe1va-galeriuxed}

\begin{enumerate}
\def\labelenumi{\arabic{enumi}.}
\tightlist
\item
  Klikněte \textbf{Galerie → Přidat album}
\item
  Vyplňte název alba
\item
  V obsahu příspěvku klikněte \textbf{Přidat média → Vytvořit galerii} →
  vyberte fotky → \textbf{Vložit galerii}
\item
  Nastavte \textbf{Obrázek příspěvku} jako náhled alba
\item
  V pravém panelu přiřaďte \textbf{Kategorii týmu}
\end{enumerate}

\paragraph{5.1.6 Správa sponzorů}\label{spruxe1va-sponzorux16f}

\begin{enumerate}
\def\labelenumi{\arabic{enumi}.}
\tightlist
\item
  Klikněte \textbf{Sponzoři → Přidat sponzora}
\item
  Vyplňte název firmy / partnera jako titulek příspěvku
\item
  V sekci \textbf{Detail sponzora} zadejte \textbf{Webové stránky} (URL,
  např. \texttt{https://firma.cz})
\item
  Nastavte \textbf{Obrázek příspěvku} -- logo sponzora
\item
  Klikněte \textbf{Publikovat}
\end{enumerate}

\paragraph{5.1.7 Správa kontaktů}\label{spruxe1va-kontaktux16f}

\begin{enumerate}
\def\labelenumi{\arabic{enumi}.}
\tightlist
\item
  Klikněte \textbf{Kontakty → Přidat kontakt}
\item
  Vyplňte jméno osoby jako titulek příspěvku (např. „Jaroslav Bejček'')
\item
  Vyplňte pole v sekci \textbf{Detail kontaktu}:
\end{enumerate}

{\def\LTcaptype{none} % do not increment counter
\begin{longtable}[]{@{}lll@{}}
\toprule\noalign{}
Pole & Popis & Příklad \\
\midrule\noalign{}
\endhead
\bottomrule\noalign{}
\endlastfoot
Funkce / Pozice & Role v klubu & \texttt{Předseda\ klubu} \\
Telefon & Kontaktní telefon & \texttt{+420\ 777\ 000\ 000} \\
E-mail & Kontaktní e-mail & \texttt{predseda@slavojmyto.cz} \\
Pořadí zobrazení & Číslo pro řazení (0 = první) & \texttt{1} \\
\end{longtable}
}

\begin{enumerate}
\def\labelenumi{\arabic{enumi}.}
\setcounter{enumi}{3}
\tightlist
\item
  Klikněte \textbf{Publikovat}
\end{enumerate}

\begin{quote}
\textbf{Tip:} Nižší číslo v poli „Pořadí zobrazení'' = kontakt se
zobrazí výše v seznamu na webu.
\end{quote}

\paragraph{5.1.8 Správa aktualit}\label{spruxe1va-aktualit}

Aktuality využívají nativní WordPress příspěvky a zobrazují se na úvodní
stránce.

\begin{enumerate}
\def\labelenumi{\arabic{enumi}.}
\tightlist
\item
  Klikněte \textbf{Příspěvky → Přidat příspěvek}
\item
  Napište titulek a text aktuality
\item
  Nastavte \textbf{Obrázek příspěvku} (náhled na homepage)
\item
  Klikněte \textbf{Publikovat}
\end{enumerate}

\subsubsection{5.2 Příručka pro návštěvníka
webu}\label{pux159uxedruux10dka-pro-nuxe1vux161tux11bvnuxedka-webu}

Web TJ Slavoj Mýto je rozdělen do přehledných sekcí přístupných z
hlavního menu. Níže je stručný popis uživatelského rozhraní.

\textbf{Navigace:} Hlavní menu v záhlaví stránky obsahuje odkazy na
Domů, Zápasy, Týmy, Galerie, Sponzoři, Kontakty a Historie. Na mobilních
zařízeních se menu sbalí do hamburger ikony (☰), kliknutím se rozbalí.

\textbf{Úvodní stránka (\texttt{/}):} Zobrazuje nejbližší nadcházející
zápas, poslední výsledky a nejnovější aktuality. Slouží jako rozcestník
k ostatním sekcím webu.

\textbf{Zápasy (\texttt{/zapasy/}):} Přehled zápasů s filtry podle týmu,
sezóny a stavu. Každý zápas zobrazuje datum, týmy, skóre (barevně --
zelená = výhra, červená = prohra) a střelce. Kliknutím na zápas se
zobrazí jeho detail.

\textbf{Týmy (\texttt{/tymy/}):} Karty týmů s fotografiemi. Kliknutím na
tým se zobrazí detail se soupiskou, realizačním týmem a nadcházejícími
zápasy.

\textbf{Galerie (\texttt{/galerie/}):} Přehled fotoalb s náhledovými
obrázky. Kliknutím na album se zobrazí fotografie s lightboxem (zvětšení
fotografie na celou obrazovku).

\textbf{Sponzoři (\texttt{/sponzori/}):} Přehled partnerů klubu s logy.
Kliknutím na logo se otevře web sponzora.

\textbf{Kontakty (\texttt{/kontakty/}):} Seznam členů vedení klubu s
funkcí, telefonem a e-mailem.

\textbf{Historie (\texttt{/historie/}):} Stránka s historií a tradicemi
klubu.

\begin{center}\rule{0.5\linewidth}{0.5pt}\end{center}

\subsection{6 Nasazení}\label{nasazenuxed}

Tato kapitola popisuje nasazení webu na produkční hosting. Doporučeným
řešením je český hosting \textbf{Wedos} s podporou PHP, MariaDB a FTP
přístupu.

\subsubsection{6.1 Nasazení na hosting
Wedos}\label{nasazenuxed-na-hosting-wedos}

\paragraph{Požadavky na hosting}\label{poux17eadavky-na-hosting}

{\def\LTcaptype{none} % do not increment counter
\begin{longtable}[]{@{}ll@{}}
\toprule\noalign{}
Požadavek & Minimální hodnota \\
\midrule\noalign{}
\endhead
\bottomrule\noalign{}
\endlastfoot
PHP & ≥ 8.0 \\
MariaDB & ≥ 10.4 \\
Diskový prostor & ≥ 1 GB \\
FTP přístup & Ano \\
\end{longtable}
}

\paragraph{Postup nasazení}\label{postup-nasazenuxed}

\begin{enumerate}
\def\labelenumi{\arabic{enumi}.}
\item
  \textbf{Vytvoření databáze} -- V klientské zóně Wedos (zákaznická
  administrace) vytvořte novou MariaDB databázi. Poznamenejte si název
  databáze, uživatelské jméno, heslo a hostname (obvykle
  \texttt{localhost} nebo adresa uvedená v administraci Wedos).
\item
  \textbf{Příprava WordPressu} -- Stáhněte WordPress z
  \href{https://wordpress.org/download/}{wordpress.org} a rozbalte na
  lokálním počítači. Do složky WordPressu zkopírujte téma a plugin:

\begin{verbatim}
web/wp-content/themes/tj-slavoj-myto/   →  wp-content/themes/tj-slavoj-myto/
web/wp-content/plugins/slavoj-custom-fields/  →  wp-content/plugins/slavoj-custom-fields/
\end{verbatim}
\item
  \textbf{Nastavení \texttt{wp-config.php}} -- Přejmenujte
  \texttt{wp-config-sample.php} na \texttt{wp-config.php} a vyplňte
  údaje databáze:

\begin{Shaded}
\begin{Highlighting}[]
\FunctionTok{define}\NormalTok{( }\StringTok{\textquotesingle{}DB\_NAME\textquotesingle{}}\OtherTok{,}     \StringTok{\textquotesingle{}nazev\_databaze\textquotesingle{}}\NormalTok{ )}\OtherTok{;}
\FunctionTok{define}\NormalTok{( }\StringTok{\textquotesingle{}DB\_USER\textquotesingle{}}\OtherTok{,}     \StringTok{\textquotesingle{}uzivatel\_databaze\textquotesingle{}}\NormalTok{ )}\OtherTok{;}
\FunctionTok{define}\NormalTok{( }\StringTok{\textquotesingle{}DB\_PASSWORD\textquotesingle{}}\OtherTok{,} \StringTok{\textquotesingle{}heslo\_databaze\textquotesingle{}}\NormalTok{ )}\OtherTok{;}
\FunctionTok{define}\NormalTok{( }\StringTok{\textquotesingle{}DB\_HOST\textquotesingle{}}\OtherTok{,}     \StringTok{\textquotesingle{}localhost\textquotesingle{}}\NormalTok{ )}\OtherTok{;}
\FunctionTok{define}\NormalTok{( }\StringTok{\textquotesingle{}DB\_CHARSET\textquotesingle{}}\OtherTok{,}  \StringTok{\textquotesingle{}utf8mb4\textquotesingle{}}\NormalTok{ )}\OtherTok{;}
\end{Highlighting}
\end{Shaded}

  Vygenerujte bezpečnostní klíče na
  \href{https://api.wordpress.org/secret-key/1.1/salt/}{api.wordpress.org/secret-key}
  a vložte je do souboru.
\item
  \textbf{Nahrání souborů přes FTP} -- Připojte se k FTP serveru Wedos
  pomocí klienta (např. FileZilla {[}doporučeno{]}). FTP údaje najdete v
  klientské zóně Wedos. Nahrajte \textbf{obsah} složky WordPressu
  (všechny soubory a podsložky) do kořenového adresáře webu
  (\texttt{www/}).
\item
  \textbf{Instalace WordPressu} -- Otevřete v prohlížeči adresu vašeho
  webu (např. \texttt{https://vasedomena.cz/}). WordPress spustí
  instalačního průvodce -- vyplňte název webu
  (\texttt{TJ\ Slavoj\ Mýto}), administrátorský účet a e-mail.
\item
  \textbf{Aktivace tématu a pluginu} -- V administraci aktivujte téma
  \textbf{TJ Slavoj Mýto} (Vzhled → Témata) a plugin \textbf{Slavoj
  Custom Fields} (Pluginy). Nastavte trvalé odkazy na \textbf{Název
  příspěvku} (Nastavení → Trvalé odkazy).
\item
  \textbf{Prvotní nastavení obsahu} -- Přejděte na \textbf{Nástroje →
  Slavoj nastavení} a klikněte „Vytvořit výchozí hodnoty taxonomií'' a
  „Vytvořit stránky webu''. Nastavte hlavní stránku (Nastavení → Čtení →
  Statická stránka → Domů) a vytvořte hlavní menu.
\item
  \textbf{SSL certifikát} -- V klientské zóně Wedos aktivujte SSL
  certifikát (Let's Encrypt -- zdarma). V WordPress admin změňte URL
  webu na \texttt{https://} (Nastavení → Obecné).
\end{enumerate}

\subsubsection{6.2 Další možnosti
nasazení}\label{dalux161uxed-moux17enosti-nasazenuxed}

Kromě Wedos lze web nasadit na jakýkoli hosting s podporou PHP ≥ 8.0 a
MySQL/MariaDB (např. Active24, Blueboard). Postup je shodný -- liší se
pouze přístup ke klientské zóně a konkrétní nastavení FTP a databáze.
Pro lokální vývoj je doporučen XAMPP {[}6{]}.

\begin{center}\rule{0.5\linewidth}{0.5pt}\end{center}

\subsection{7 Závěr}\label{zuxe1vux11br}

Cílem této maturitní práce bylo vytvořit moderní webovou prezentaci pro
fotbalový klub TJ Slavoj Mýto, která umožní správci klubu spravovat
obsah bez znalosti programování. Tento cíl byl splněn.

Výsledný web je postaven na systému WordPress s vlastním tématem
\texttt{tj-slavoj-myto} a pluginem \texttt{slavoj-custom-fields}. Místo
závislosti na externích pluginech byla klíčová funkcionalita
implementována vlastním kódem -- registrace šesti typů obsahu, čtyř
taxonomií, administrační meta boxy, vlastní sloupce, filtry a
uživatelská role „Správce obsahu''.

Web je plně responzivní díky frameworku Bootstrap 5 a mobile-first
přístupu ke stylování. Kód dodržuje bezpečnostní standardy WordPressu
(sanitizace, escapování, nonce) a minimalizuje použití JavaScriptu.

\subsubsection{Splnění bodů
zadání}\label{splnux11bnuxed-bodux16f-zaduxe1nuxed}

{\def\LTcaptype{none} % do not increment counter
\begin{longtable}[]{@{}
  >{\raggedright\arraybackslash}p{(\linewidth - 2\tabcolsep) * \real{0.5000}}
  >{\raggedright\arraybackslash}p{(\linewidth - 2\tabcolsep) * \real{0.5000}}@{}}
\toprule\noalign{}
\begin{minipage}[b]{\linewidth}\raggedright
Bod zadání
\end{minipage} & \begin{minipage}[b]{\linewidth}\raggedright
Stav
\end{minipage} \\
\midrule\noalign{}
\endhead
\bottomrule\noalign{}
\endlastfoot
Figma design (desktop + mobilní) & Ano Vytvořen \\
WordPress téma s Bootstrap 5 & Ano Implementováno \\
Custom post types (zápasy, týmy, hráči\ldots) & Ano 6 CPT + 4
taxonomie \\
Filtrování podle sezóny a týmu & Ano GET filtry na archivních
stránkách \\
Galerie fotografií & Ano CPT s nativní WordPress galerií + lightbox \\
Optimalizace výkonu & Ano Lazy loading, CDN, minimální JS \\
Uživatelské role & Ano Administrátor + Správce obsahu \\
\end{longtable}
}

\subsubsection{Návrhy na
rozšíření}\label{nuxe1vrhy-na-rozux161uxedux159enuxed}

\begin{enumerate}
\def\labelenumi{\arabic{enumi}.}
\tightlist
\item
  \textbf{Kalendářní modul} -- vizuální kalendář zápasů (doporučen
  plugin The Events Calendar {[}17{]})
\item
  \textbf{SEO optimalizace} -- integrace pluginu Rank Math {[}18{]} nebo
  Yoast SEO {[}19{]}
\item
  \textbf{Cache a komprese} -- nasazení WP Super Cache {[}20{]} a Smush
  {[}21{]} pro produkční prostředí
\item
  \textbf{Vícejazyčnost} -- příprava pro anglickou verzi webu
\item
  \textbf{Statistiky hráčů} -- rozšíření dat hráčů o góly, asistence a
  zápasy
\end{enumerate}

\begin{center}\rule{0.5\linewidth}{0.5pt}\end{center}

\subsection{Seznam použité
literatury}\label{seznam-pouux17eituxe9-literatury}

{[}1{]} BEJČEK, Lukáš. \emph{Zadání maturitní práce: Webová prezentace
fotbalového klubu s databází zápasů a správou obsahu ve WordPressu}
{[}dokument{]}. 2025. Soubor \texttt{ZADANI-MATURITNI-PRACE.md} v kořeni
repozitáře.

{[}2{]} WORDPRESS FOUNDATION. \emph{WordPress} {[}software{]}. Verze
6.x. San Francisco (CA): WordPress Foundation, 2024 {[}cit.
2026-03-27{]}. Dostupné z: https://wordpress.org

{[}3{]} WORDPRESS FOUNDATION. \emph{WordPress Developer Resources}
{[}online{]}. San Francisco (CA): WordPress Foundation, 2024 {[}cit.
2026-03-27{]}. Dostupné z: https://developer.wordpress.org

{[}4{]} OTTO, Mark a Jacob THORNTON. \emph{Bootstrap} {[}software{]}.
Verze 5.3.3. {[}b.m.{]}: Bootstrap Team, 2024 {[}cit. 2026-03-27{]}.
Dostupné z: https://getbootstrap.com

{[}5{]} JSDELIVR. \emph{jsDelivr -- free CDN for open source}
{[}online{]}. 2024 {[}cit. 2026-03-27{]}. Dostupné z:
https://www.jsdelivr.com

{[}6{]} APACHE FRIENDS. \emph{XAMPP} {[}software{]}. {[}b.m.{]}: Apache
Friends, 2024 {[}cit. 2026-03-27{]}. Dostupné z:
https://www.apachefriends.org

{[}7{]} FIGMA, INC. \emph{Figma} {[}software{]}. San Francisco (CA):
Figma, 2024 {[}cit. 2026-03-27{]}. Dostupné z: https://www.figma.com

{[}8{]} FLAVOR, WebDevStudios. Custom Post Type UI. In: \emph{WordPress
Plugin Directory} {[}online{]}. San Francisco (CA): WordPress
Foundation, 2024 {[}cit. 2026-03-27{]}. Dostupné z:
https://wordpress.org/plugins/custom-post-type-ui/

{[}9{]} WP ENGINE. Advanced Custom Fields. In: \emph{WordPress Plugin
Directory} {[}online{]}. San Francisco (CA): WordPress Foundation, 2024
{[}cit. 2026-03-27{]}. Dostupné z:
https://wordpress.org/plugins/advanced-custom-fields/

{[}10{]} FACETWP, LLC. \emph{FacetWP -- Advanced Filtering for
WordPress} {[}online{]}. 2024 {[}cit. 2026-03-27{]}. Dostupné z:
https://facetwp.com

{[}11{]} FLAVOR, Vladimir Garagulja. User Role Editor. In:
\emph{WordPress Plugin Directory} {[}online{]}. San Francisco (CA):
WordPress Foundation, 2024 {[}cit. 2026-03-27{]}. Dostupné z:
https://wordpress.org/plugins/user-role-editor/

{[}12{]} WORLD WIDE WEB CONSORTIUM (W3C). \emph{HTML Living Standard:
Lazy loading} {[}online{]}. 2024 {[}cit. 2026-03-27{]}. Dostupné z:
https://html.spec.whatwg.org/multipage/urls-and-fetching.html\#lazy-loading-attributes

{[}13{]} WORDPRESS FOUNDATION. \emph{Theme Developer Handbook}
{[}online{]}. San Francisco (CA): WordPress Foundation, 2024 {[}cit.
2026-03-27{]}. Dostupné z: https://developer.wordpress.org/themes/

{[}14{]} BEJČEK, Lukáš. \emph{JavaScript v projektu TJ Slavoj Mýto}
{[}dokument{]}. 2026. Soubor \texttt{docs/pracovni/12-javascript.md} v
repozitáři.

{[}15{]} SOFTWARE FREEDOM CONSERVANCY. \emph{Git} {[}software{]}. Verze
2.x. {[}b.m.{]}: Software Freedom Conservancy, 2024 {[}cit.
2026-03-27{]}. Dostupné z: https://git-scm.com

{[}16{]} DOCKER, INC. \emph{Docker} {[}software{]}. San Francisco (CA):
Docker, 2024 {[}cit. 2026-03-27{]}. Dostupné z: https://www.docker.com

{[}17{]} THE EVENTS CALENDAR. The Events Calendar. In: \emph{WordPress
Plugin Directory} {[}online{]}. San Francisco (CA): WordPress
Foundation, 2024 {[}cit. 2026-03-27{]}. Dostupné z:
https://wordpress.org/plugins/the-events-calendar/

{[}18{]} FLAVOR, Starter Sites. Starter Templates. Rank Math SEO. In:
\emph{WordPress Plugin Directory} {[}online{]}. San Francisco (CA):
WordPress Foundation, 2024 {[}cit. 2026-03-27{]}. Dostupné z:
https://wordpress.org/plugins/seo-by-rank-math/

{[}19{]} YOAST BV. Yoast SEO. In: \emph{WordPress Plugin Directory}
{[}online{]}. San Francisco (CA): WordPress Foundation, 2024 {[}cit.
2026-03-27{]}. Dostupné z: https://wordpress.org/plugins/wordpress-seo/

{[}20{]} FLAVOR, Developer. WP Super Cache. In: \emph{WordPress Plugin
Directory} {[}online{]}. San Francisco (CA): WordPress Foundation, 2024
{[}cit. 2026-03-27{]}. Dostupné z:
https://wordpress.org/plugins/wp-super-cache/

{[}21{]} FLAVOR, Developer. Smush. In: \emph{WordPress Plugin Directory}
{[}online{]}. San Francisco (CA): WordPress Foundation, 2024 {[}cit.
2026-03-27{]}. Dostupné z: https://wordpress.org/plugins/wp-smushit/

\begin{center}\rule{0.5\linewidth}{0.5pt}\end{center}

\subsection{Příloha A -- Kopie zadání
práce}\label{pux159uxedloha-a-kopie-zaduxe1nuxed-pruxe1ce}

\textbf{Jméno a příjmení:} Lukáš Bejček \textbf{Studijní obor:}
Informační technologie \textbf{Název tématu:} Webová prezentace
fotbalového klubu s databází zápasů a správou obsahu ve WordPressu

\subsubsection{Cíl práce}\label{cuxedl-pruxe1ce}

Cílem práce je vytvořit moderní informační web pro fotbalový klub, který
bude přehledně zobrazovat zápasy, týmy, historii, partnery a aktuální
informace o dění v klubu. Web bude postaven na vlastním grafickém návrhu
vytvořeném ve Figmě a nasazeném na systému WordPress s využitím vlastní
šablony založené na Bootstrapu 5. Web bude připraven pro správu obsahu
vedením klubu nebo pověřeným administrátorem.

\subsubsection{Zásady pro
vypracování}\label{zuxe1sady-pro-vypracovuxe1nuxed}

\begin{enumerate}
\def\labelenumi{\arabic{enumi}.}
\item
  Vytvoření grafického návrhu ve Figmě. Návrh musí obsahovat samostatné
  varianty pro desktop i mobilní zařízení, včetně hlavní stránky i
  podstránek. Grafický návrh musí zahrnovat také vizuální podobu
  použitých pluginů a komponent, aby zapadaly do celkového designu
  stránky.
\item
  Tvorba šablony v Bootstrapu 5 a nasazení ve WordPressu. Použití
  souborů header.php, footer.php, single.php, archive.php,
  functions.php, styles.css, atd.
\item
  Zápasy a týmy jako vlastní typy obsahu

  \begin{enumerate}
  \def\labelenumii{\alph{enumii}.}
  \tightlist
  \item
    Vytvoření typu příspěvku „Zápas'' s poli: datum, čas, tým, soupeř,
    výsledek, střelci, sezóna, status zápasu
  \item
    Vytvoření typu příspěvku „Tým'' s poli: název, sezóna, trenér,
    asistenti, seznam hráčů
  \item
    Pluginy: Custom Post Type UI, Advanced Custom Fields, Pods
    (volitelné)
  \end{enumerate}
\item
  Filtrování a výpis podle sezóny nebo týmů

  \begin{enumerate}
  \def\labelenumii{\alph{enumii}.}
  \tightlist
  \item
    Možnost filtrování zápasů a týmů podle ročníku, stavu, nebo
    kategorie
  \item
    Pluginy: FacetWP, Search \& Filter Pro
  \end{enumerate}
\item
  Kalendářový modul

  \begin{enumerate}
  \def\labelenumii{\alph{enumii}.}
  \tightlist
  \item
    Implementace kalendáře, který umožní vizuální zobrazení
    nadcházejících i odehraných zápasů
  \item
    Možnost filtrování podle týmů, kategorií nebo sezón
  \item
    Pluginy: The Events Calendar, Event Organiser, Simple Calendar
  \end{enumerate}
\item
  Galerie

  \begin{enumerate}
  \def\labelenumii{\alph{enumii}.}
  \tightlist
  \item
    Vytvoření galerie s možností seskupení podle názvu nebo roku
  \item
    Pluginy: Modula, Envira Gallery, NextGEN Gallery
  \end{enumerate}
\item
  Optimalizace výkonu a načítání

  \begin{enumerate}
  \def\labelenumii{\alph{enumii}.}
  \tightlist
  \item
    Cache (mezipaměť) -- ukládání stránek do mezipaměti pro rychlejší
    načítání
  \item
    Lazy loading -- odložené načítání obrázků a externího obsahu až při
    posunu stránky
  \item
    Minifikace -- zmenšení velikosti CSS a JS souborů
  \item
    Optimalizace obrázků -- komprese a správný formát (např. WebP)
  \item
    Kombinace CSS/JS -- snížení počtu HTTP požadavků
  \item
    Pluginy: WP Super Cache, LiteSpeed Cache, Smush
  \end{enumerate}
\item
  SEO a analytika

  \begin{enumerate}
  \def\labelenumii{\alph{enumii}.}
  \tightlist
  \item
    Pluginy: Yoast SEO, Rank Math, propojení s Google Analytics
  \end{enumerate}
\item
  Hosting a zabezpečení

  \begin{enumerate}
  \def\labelenumii{\alph{enumii}.}
  \tightlist
  \item
    Hosting, HTTPS certifikát, ochrana proti útokům a zálohování
  \end{enumerate}
\item
  Uživatelské role a oprávnění

  \begin{enumerate}
  \def\labelenumii{\alph{enumii}.}
  \tightlist
  \item
    Administrátor, přispěvatel -- různá oprávnění pro správu obsahu
  \item
    Pluginy: User Role Editor, Members
  \end{enumerate}
\item
  Odevzdávané výstupy

  \begin{enumerate}
  \def\labelenumii{\alph{enumii}.}
  \tightlist
  \item
    Podrobná dokumentace všech použitých pluginů
  \item
    Zdrojový kód WordPress šablony a popis všech jejích úprav
  \item
    Funkční web (na hostingu nebo jako lokalizovaný projekt s návodem ke
    spuštění a veškerými přístupovými údaji)
  \end{enumerate}
\end{enumerate}

\textbf{Rozsah dokumentace:} 15--50 stran \textbf{Forma MP:}
elektronická

\subsection{Příloha B -- Přístupové
údaje}\label{pux159uxedloha-b-pux159uxedstupovuxe9-uxfadaje}

\begin{quote}
\textbf{Upozornění:} Tyto údaje jsou určeny výhradně pro lokální
testování. Při nasazení na produkční server je nutné zvolit silná a
unikátní hesla.
\end{quote}

\subsubsection{WordPress administrace}\label{wordpress-administrace}

\textbf{URL:} \texttt{http://localhost/fotbal\_club/wp-admin/}

{\def\LTcaptype{none} % do not increment counter
\begin{longtable}[]{@{}
  >{\raggedright\arraybackslash}p{(\linewidth - 6\tabcolsep) * \real{0.1395}}
  >{\raggedright\arraybackslash}p{(\linewidth - 6\tabcolsep) * \real{0.4419}}
  >{\raggedright\arraybackslash}p{(\linewidth - 6\tabcolsep) * \real{0.1628}}
  >{\raggedright\arraybackslash}p{(\linewidth - 6\tabcolsep) * \real{0.2558}}@{}}
\toprule\noalign{}
\begin{minipage}[b]{\linewidth}\raggedright
Role
\end{minipage} & \begin{minipage}[b]{\linewidth}\raggedright
Uživatelské jméno
\end{minipage} & \begin{minipage}[b]{\linewidth}\raggedright
Heslo
\end{minipage} & \begin{minipage}[b]{\linewidth}\raggedright
Oprávnění
\end{minipage} \\
\midrule\noalign{}
\endhead
\bottomrule\noalign{}
\endlastfoot
Administrátor & \texttt{admin} & \texttt{admin123} & Plný přístup k
celému WordPressu \\
Správce obsahu & \texttt{redaktor} & \texttt{redaktor123} & Editace CPT,
nahrávání médií, bez nastavení WP a pluginů \\
\end{longtable}
}

\textbf{Administrátor} má plný přístup ke všem funkcím WordPressu včetně
nastavení, pluginů, témat a všech typů obsahu.

\textbf{Správce obsahu} může: - Přidávat, editovat a mazat zápasy, týmy,
hráče, galerie, sponzory a kontakty - Nahrávat obrázky a média -
Spravovat kategorie a taxonomie

Správce obsahu \textbf{nemůže}: - Měnit nastavení WordPressu -
Instalovat nebo mazat pluginy a témata - Spravovat uživatelské účty

\subsubsection{Databáze (MySQL -- XAMPP)}\label{databuxe1ze-mysql-xampp}

{\def\LTcaptype{none} % do not increment counter
\begin{longtable}[]{@{}ll@{}}
\toprule\noalign{}
Parametr & Hodnota \\
\midrule\noalign{}
\endhead
\bottomrule\noalign{}
\endlastfoot
Server & \texttt{localhost} \\
Port & \texttt{3306} \\
Uživatel & \texttt{root} \\
Heslo & \emph{(prázdné -- výchozí XAMPP)} \\
Databáze & \texttt{slavoj\_myto} \\
phpMyAdmin & \texttt{http://localhost/phpmyadmin/} \\
\end{longtable}
}
